\documentclass[12pt,a4paper]{article}
\usepackage[utf8]{inputenc}
\usepackage[T1]{fontenc}
\usepackage[french]{babel}
\usepackage{geometry}
\usepackage{graphicx}
\usepackage{hyperref}
\usepackage{listings}
\usepackage{xcolor}
\usepackage{enumitem}
\usepackage{titlesec}
\usepackage{fancyhdr}
\usepackage{fontawesome5}
\usepackage{tikz}

% Configuration de la page
\geometry{margin=2.5cm}
\pagestyle{fancy}
\fancyhf{}
\fancyhead[L]{\leftmark}
\fancyhead[R]{Robot Hotel - Cahier des Charges}
\fancyfoot[C]{\thepage}

% Configuration des liens
\hypersetup{
    colorlinks=true,
    linkcolor=blue,
    filecolor=magenta,      
    urlcolor=cyan,
    pdftitle={Cahier des Charges - Robot Hotel},
}

% Configuration des listes
\setlist[itemize]{leftmargin=*,itemsep=0.1cm}
\setlist[enumerate]{leftmargin=*,itemsep=0.1cm}

% Configuration des titres
\titleformat{\section}
{\Large\bfseries\color{blue!70!black}}
{}{0em}{}[\titlerule]

\titleformat{\subsection}
{\large\bfseries\color{blue!60!black}}
{}{0em}{}

% Couleurs personnalisées
\definecolor{primaryblue}{RGB}{41,128,185}
\definecolor{secondaryblue}{RGB}{52,152,219}
\definecolor{accentorange}{RGB}{231,76,60}
\definecolor{successgreen}{RGB}{39,174,96}

% Définition d'environnement pour les blocs
\newenvironment{block}[1]{\begin{center}\begin{minipage}{0.9\textwidth}\textbf{#1}\\}{\end{minipage}\end{center}}

\newenvironment{alertblock}[1]{\begin{center}\begin{minipage}{0.9\textwidth}\color{accentorange}\textbf{#1}\\}{\end{minipage}\end{center}}

% Informations du document
\title{
    \vspace{2cm}
    {\Huge\bfseries Cahier des Charges}\\
    \vspace{0.5cm}
    {\Large Projet Robot Hotel}\\
    \vspace{0.3cm}
    {\large Système Intelligent d'Accueil et de Gestion Hôtelière}
    \vspace{1.5cm}
}

\author{
    christ \& farouk\\
    Université d'Avignon
}

\date{\today}

\begin{document}

\maketitle
\thispagestyle{empty}
\newpage

\tableofcontents
\newpage

% ============================================================================
% SECTION 1: CONTEXTE ET OBJECTIFS
% ============================================================================

\section{Contexte et Objectifs}

\subsection{Contexte du Projet}

\begin{block}{Vision}
Système intelligent d'accueil et de gestion hôtelière basé sur robot humanoïde Pepper
\end{block}

\begin{description}
\item[Objectifs principaux :]
\begin{itemize}
    \item Améliorer l'expérience d'accueil
    \item Automatiser les tâches de réception
    \item Service d'information 24/7
    \item Gestion automatisée des réservations
\end{itemize}

\item[Besoins identifiés :]
\begin{itemize}
    \item Personnalisation de l'accueil
    \item Accessibilité multilingue
    \item Efficacité opérationnelle
    \item Intégration technologique
\end{itemize}
\end{description}

\subsection{Objectifs Fonctionnels}

\begin{block}{Fonctionnalités clés}
\begin{itemize}
    \item[\textcolor{primaryblue}{\faUser}] \textbf{Reconnaissance faciale} pour identification clients
    \item[\textcolor{primaryblue}{\faComments}] \textbf{Chatbot intelligent} (Llama 3.1) avec RAG
    \item[\textcolor{primaryblue}{\faMicrophone}] \textbf{Transcription audio} (Whisper)
    \item[\textcolor{primaryblue}{\faDatabase}] \textbf{API REST complète} (Flask)
    \item[\textcolor{primaryblue}{\faDesktop}] \textbf{Interface admin} (Streamlit)
    \item[\textcolor{primaryblue}{\faBed}] \textbf{Gestion réservations} automatisée
\end{itemize}
\end{block}

\begin{alertblock}{Objectifs techniques}
Architecture modulaire, performances temps réel, compatibilité Pepper, documentation API
\end{alertblock}

% ============================================================================
% SECTION 2: PÉRIMÈTRE
% ============================================================================

\section{Périmètre du Projet}

\subsection{Fonctionnalités Principales}

\subsubsection{Reconnaissance Faciale}
\begin{itemize}
    \item Enregistrement personnes
    \item Détection temps réel
    \item Streaming vidéo
    \item Gestion base de données
\end{itemize}

\subsubsection{Chat Intelligent}
\begin{itemize}
    \item Chatbot Llama 3.1
    \item Contexte RAG
    \item Connaissance temporelle
    \item Support multilingue
\end{itemize}

\subsubsection{Audio}
\begin{itemize}
    \item Transcription Whisper
    \item Enregistrement vocal
    \item Détection silence
    \item Extraction informations
\end{itemize}

\subsubsection{Réservations}
\begin{itemize}
    \item Création depuis audio
    \item Extraction intelligente
    \item QR codes automatiques
    \item Vérification QR codes
\end{itemize}

% ============================================================================
% SECTION 3: TECHNOLOGIES
% ============================================================================

\section{Technologies Utilisées}

\subsection{Stack Technologique}

\subsubsection{Frameworks}
\begin{itemize}
    \item Flask (API REST)
    \item FastAPI (Réservations)
    \item Streamlit (Interface)
    \item Uvicorn (Serveur ASGI)
\end{itemize}

\subsubsection{IA / ML}
\begin{itemize}
    \item Llama 3.1-8B-Instruct
    \item Whisper Medium
    \item YOLOv8
    \item face\_recognition
    \item OpenCV
\end{itemize}

\subsubsection{RAG}
\begin{itemize}
    \item LlamaIndex
    \item HuggingFace Embeddings
    \item Vector Store
\end{itemize}

\subsubsection{Bases de Données}
\begin{itemize}
    \item Pickle (Visages)
    \item SQLite (Réservations)
    \item SQLAlchemy (ORM)
\end{itemize}

\subsubsection{Matériel}
\begin{itemize}
    \item Robot Pepper
    \item Webcam
    \item Microphone
\end{itemize}

% ============================================================================
% SECTION 4: FONCTIONNALITÉS IMPLÉMENTÉES
% ============================================================================

\section{Fonctionnalités Implémentées}

\subsection{Architecture Multi-Services}

\begin{block}{Système composé de 3 services principaux}
\begin{enumerate}
    \item \textbf{API Flask} (port 5000) - Serveur principal avec Swagger
    \item \textbf{API FastAPI} (port 8000) - Gestion réservations
    \item \textbf{Interface Streamlit} (port 8501) - Administration
\end{enumerate}
\end{block}

\begin{block}{Automatisation}
Script Bash unifié (\texttt{start\_all.sh}) pour démarrer tous les services avec leurs environnements virtuels
\end{block}

\subsection{Reconnaissance Faciale}

\subsubsection{Technologies}
\begin{itemize}
    \item YOLOv8 (détection)
    \item face\_recognition (reconnaissance)
    \item Base pickle
\end{itemize}

\subsubsection{Fonctionnalités}
\begin{itemize}
    \item Enregistrement via upload
    \item Détection temps réel
    \item Streaming vidéo
    \item Gestion complète BDD
\end{itemize}

\subsubsection{Interface Admin}
\begin{itemize}
    \item Page Streamlit dédiée
    \item Preview images
    \item Liste et suppression
    \item Test en temps réel
\end{itemize}

\subsection{Chat Intelligent avec RAG}

\begin{block}{Chatbot Llama 3.1-8B-Instruct}
\begin{itemize}
    \item Connexion Hugging Face
    \item Conversations contextuelles
    \item Réponses basées sur documentation hôtel
\end{itemize}
\end{block}

\begin{block}{Système RAG}
\begin{itemize}
    \item Indexation vectorielle (LlamaIndex)
    \item Embeddings HuggingFace
    \item Recherche sémantique
\end{itemize}
\end{block}

\begin{block}{Contexte Temporel}
Intégration automatique date/heure pour réponses temporelles précises
\end{block}

\subsection{Gestion de Réservations}

\subsubsection{Création depuis Audio}
\begin{itemize}
    \item Transcription Whisper
    \item Extraction Llama 3.1
    \item Conversion dates relatives
\end{itemize}

\subsubsection{API FastAPI}
\begin{itemize}
    \item Création clients
    \item Création réservations
    \item Vérification QR codes
    \item Confirmation arrivée
\end{itemize}

\subsubsection{QR Codes}
\begin{itemize}
    \item Génération automatique
    \item Stockage PNG
    \item Vérification scan
    \item Interface dédiée
\end{itemize}

\subsubsection{Notifications Email}
\begin{itemize}
    \item Envoi automatique
    \item Email HTML formaté
    \item QR code en pièce jointe
    \item Configuration SMTP
\end{itemize}

\subsection{Interface Streamlit - 5 Onglets}

\begin{block}{Organisation}
\begin{enumerate}
    \item \textcolor{primaryblue}{\faUser} \textbf{Gestion Visages} - Enregistrement, liste, suppression
    \item \textcolor{primaryblue}{\faVideo} \textbf{Streaming Vidéo} - Contrôle détection faciale
    \item \textcolor{primaryblue}{\faComments} \textbf{Chat et Audio} - Chat, transcription, enregistrement
    \item \textcolor{primaryblue}{\faBed} \textbf{Réservations} - Création, vérification QR codes
    \item \textcolor{primaryblue}{\faEnvelope} \textbf{Configuration Email} - Paramètres SMTP
\end{enumerate}
\end{block}

\begin{block}{Fonctionnalités}
Test connexion API, gestion erreurs, feedback utilisateur, logs détaillés
\end{block}

\subsection{Intégration Robot Pepper}

\subsubsection{Script Principal}
\begin{itemize}
    \item Interaction complète
    \item Enregistrement audio
    \item Transcription + Chat
    \item Réponse vocale
\end{itemize}

\subsubsection{Fonctionnalités}
\begin{itemize}
    \item Création réservations
    \item Scan QR codes
    \item Interface tablette
\end{itemize}

\subsubsection{Workflow}
\begin{enumerate}
    \item Reconnaissance
    \item Interaction vocale
    \item Traitement IA
    \item Réponse robot
\end{enumerate}

% ============================================================================
% SECTION 5: SCÉNARIOS DE TESTS
% ============================================================================

\section{Scénarios de Tests}

\subsection{Scénario 1 : Reconnaissance}

\begin{block}{Scénario 1 : Reconnaissance}
\begin{itemize}
    \item Enregistrement visage
    \item Détection robot/caméra
    \item Accueil personnalisé
    \item \textbf{Critère} : 80\% précision, < 3s
\end{itemize}
\end{block}

\subsection{Scénario 2 : Chat}

\begin{block}{Scénario 2 : Chat}
\begin{itemize}
    \item Questions services/équipements
    \item Questions temporelles
    \item Questions multilingues
    \item \textbf{Critère} : > 85\% pertinence, < 5s
\end{itemize}
\end{block}

\subsection{Scénario 3 : Réservation Vocale}

\begin{block}{Scénario 3 : Réservation Vocale}
\begin{itemize}
    \item Enregistrement demande
    \item Transcription automatique
    \item Extraction informations
    \item Génération QR code
    \item \textbf{Critère} : > 90\% transcription, > 85\% extraction
\end{itemize}
\end{block}

\subsection{Scénario 4 : Administration}

\begin{block}{Scénario 4 : Administration}
\begin{itemize}
    \item Navigation interface
    \item Gestion visages
    \item Tests fonctionnalités
    \item \textbf{Critère} : Navigation intuitive, < 2min/opération
\end{itemize}
\end{block}

% ============================================================================
% SECTION 6: CRITÈRES DE RÉUSSITE
% ============================================================================

\section{Critères de Réussite}

% Section laissée vide comme dans la présentation

% ============================================================================
% SECTION 7: AMÉLIORATIONS FUTURES
% ============================================================================

\section{Améliorations Futures}

\subsection{Améliorations Futures}

\begin{block}{Priorités}
\begin{itemize}
    \item[\textcolor{successgreen}{\faCheckSquare}] Dashboard analytics
    \item[\textcolor{successgreen}{\faCheckSquare}] Authentification multi-utilisateurs
    \item[\textcolor{successgreen}{\faCheckSquare}] Amélioration précision reconnaissance
    \item[\textcolor{successgreen}{\faCheckSquare}] Dockerisation complète
    \item[\textcolor{successgreen}{\faCheckSquare}] Améliorer le temps de réponse
    \item[\textcolor{successgreen}{\faCheckSquare}] Amélioration de l'interaction du robot avec l'humain.
\end{itemize}
\end{block}

% ============================================================================
% SECTION 8: CONTRAINTES ET RISQUES
% ============================================================================

\section{Contraintes et Risques}

\subsection{Contraintes Techniques}

\subsubsection{Contraintes Matérielles}
\begin{itemize}
    \item Disponibilité Pepper
    \item Ressources computationnelles
    \item Stockage modèles ML
    \item Connexion réseau
\end{itemize}

\subsubsection{Contraintes Logicielles}
\begin{itemize}
    \item Python 3.12+
    \item Dépendances multiples
    \item Environnements virtuels
    \item Compatibilité OS
\end{itemize}

\begin{alertblock}{Mesures}
Tests approfondis, documentation complète, logs détaillés, mode dégradé
\end{alertblock}

\subsection{Risques et Mitigation}

\begin{block}{Risques Techniques}
\begin{itemize}
    \item Performance modèles IA → Tests approfondis
    \item Intégration Pepper → Documentation complète
    \item Bugs critiques → Système logs détaillé
\end{itemize}
\end{block}

\begin{block}{Risques Utilisateur}
\begin{itemize}
    \item Acceptabilité → Tests utilisateurs précoces
    \item Complexité perçue → Design simple
    \item Fiabilité → Messages erreurs clairs
\end{itemize}
\end{block}

% ============================================================================
% SECTION 9: PLANNING
% ============================================================================

\section{Planning}

% Section laissée vide comme dans la présentation

% ============================================================================
% CONCLUSION
% ============================================================================

\section{Conclusion}

\subsection{Récapitulatif}

\begin{block}{Système complet implémenté}
\begin{itemize}
    \item[\textcolor{successgreen}{\faCheckCircle}] Reconnaissance faciale opérationnelle
    \item[\textcolor{successgreen}{\faCheckCircle}] Chatbot intelligent avec RAG
    \item[\textcolor{successgreen}{\faCheckCircle}] Transcription et enregistrement audio
    \item[\textcolor{successgreen}{\faCheckCircle}] Gestion réservations automatique
    \item[\textcolor{successgreen}{\faCheckCircle}] Interface admin complète
    \item[\textcolor{successgreen}{\faCheckCircle}] Intégration Robot Pepper
    \item[\textcolor{successgreen}{\faCheckCircle}] Notifications email avec QR codes
\end{itemize}
\end{block}

\begin{block}{Objectifs}
Améliorer l'expérience client tout en validant l'acceptabilité des technologies innovantes
\end{block}

\vspace{2cm}

\noindent\textit{Document version 1.0 - \today}

\end{document}
